% Academic CV for Xiao Chen.
% Updated from content/about.md, content/site.yml, and the previous CV.
% Template originally adapted from Jake Gutierrez / sb2nov (MIT license).
% Compile with: tectonic public/file/XiaoChen_CV.tex

\documentclass[letterpaper,11pt]{article}
\usepackage[T1]{fontenc}
\usepackage[utf8]{inputenc}
\usepackage{lmodern}
\usepackage[margin=0.65in,footskip=0.25in]{geometry}
\usepackage{titlesec}
\usepackage{enumitem}
\usepackage{tabularx}
\usepackage{array}
\usepackage{microtype}
\usepackage{needspace}
\usepackage{fancyhdr}
\usepackage[hidelinks]{hyperref}

\hypersetup{
  pdftitle={Xiao Chen - Curriculum Vitae},
  pdfauthor={Xiao Chen},
  pdfsubject={Programming languages, migratory typing, effect handlers, and formal verification}
}
\urlstyle{same}
\setlength{\parindent}{0pt}
\setlength{\parskip}{0pt}
\setlength{\tabcolsep}{0pt}
\setlength{\emergencystretch}{1.5em}
\hyphenpenalty=10000
\exhyphenpenalty=10000
\raggedbottom

\pagestyle{fancy}
\fancyhf{}
\renewcommand{\headrulewidth}{0pt}
\renewcommand{\footrulewidth}{0pt}

\titleformat{\section}{\large\bfseries\sffamily}{}{0pt}{}[\titlerule]
\titlespacing*{\section}{0pt}{10pt}{5pt}
\newenvironment{cvitems}{
  \begin{itemize}[leftmargin=1.2em,itemsep=2pt,topsep=3pt,parsep=0pt,partopsep=0pt]
}{\end{itemize}}
\newcommand{\cvheading}[2]{%
  \Needspace{4\baselineskip}%
  \begin{tabularx}{\linewidth}{@{}>{\raggedright\arraybackslash}X@{\hspace{1em}}r@{}}%
    \textbf{#1} & #2
  \end{tabularx}\par
}
\newcommand{\cvdetail}[1]{#1\par}
\newcommand{\cvproject}[2]{%
  \Needspace{6\baselineskip}%
  \textbf{#1}\par
  \textit{#2}\par
}
\newcommand{\entrygap}{\vspace{5pt}}

\newcommand{\vsp}{\vspace{6pt}}

\begin{document}
\small

\begin{center}
  {\LARGE\bfseries Xiao Chen}\par
  \vspace{4pt}
  Ph.D. Student in Computer Science, Indiana University Bloomington\par
  \vspace{3pt}
  \textbf{xc56@iu.edu}
  \enspace | \enspace
  \textbf{chenxiao2402.github.io}
  \enspace | \enspace
  \textbf{GitHub: chenxiao2402}
\end{center}

%\vspace{-5pt}

\textbf{Research interests:} Programming languages; gradual and migratory typing; operational semantics; algebraic effect handlers; formal verification.\par
\vspace{3pt}
Seeking Ph.D. opportunities outside the United States for 2027 entry.

\vsp

\section{Education}
\newlength{\educationdatewidth}
\settowidth{\educationdatewidth}{Sep. 2022 -- Jun. 2024}
\cvheading{Indiana University Bloomington}{\makebox[\educationdatewidth][l]{Sep. 2024 -- Present}}
\cvdetail{Ph.D. in Computer Science \enspace | \enspace Advisor: Prof. Sam Tobin-Hochstadt}
\cvdetail{Currently on leave due to prolonged U.S. visa processing delays.}
\entrygap
\cvheading{Northwestern University}{Sep. 2022 -- Jun. 2024}
\cvdetail{M.S. in Computer Science \enspace | \enspace GPA: 3.957/4.00}
\cvdetail{Research with Profs. Christos Dimoulas and Simone Campanoni.}
\entrygap
\cvheading{Nanjing University}{Sep. 2020 -- Jun. 2022}
\cvdetail{M.Eng. in Electronic Information \enspace | \enspace GPA: 4.58/5.00 (1\%)}
\entrygap
\cvheading{Nanjing University}{Sep. 2016 -- Jun. 2020}
\cvdetail{B.Eng. in Software Engineering \enspace | \enspace GPA: 4.28/5.00 (20\%)}

\vsp


\section{Ongoing Research Projects}
\cvproject{Type as Obligations (TAO)}{Formalized proofs in Rocq complete; manuscript in preparation.}
\begin{cvitems}
  \item Developed an executable operational framework that separates the type obligations accumulated during execution from the checks performed by a runtime enforcement strategy.
  \item Mechanically derived a profiler semantics that tracks typing histories and boundary provenance through higher-order interactions, without manually instrumenting individual programs.
  \item Proved a bidirectional correspondence between ordinary and profiler evaluation: well-typed, uniquely labeled programs terminate with the same resolved values, despite extra bookkeeping and different location allocations.
  \item Defined obligation soundness and proved that it implies type soundness, but not conversely. Used violated obligations as a common reference for blame soundness and completeness.
  \item Compared Natural, Threesomes, Forgetful, and Transient using uniform criteria, with formalized proofs and counterexamples in \textbf{Rocq/Coq} and an executable \textbf{Racket/Redex} model.
\end{cvitems}

\vsp

\cvproject{Effect Typed Racket (ETR)}{Prototype implemented; formal semantics and proofs in progress.}

\begin{cvitems}
  \item Extended \textbf{Typed Racket} with algebraic effect handlers and effect typing, carrying effect information through higher-order calls, polymorphic instantiation, and module imports.
  \item Implemented runtime contracts for escaping-effect specifications across typed/untyped boundaries. Distinguished forwarded requests and resumed computations from effects of the handler implementation.
  \item Supported effect identity across modules and explicit effect-row instantiation. Current contracts check concrete effect allowances; relational effect-row parametricity remains ongoing work.
  \item Developing a high-level calculus and a lower-level model based on delimited continuations and contracts, aiming to establish type and effect soundness and a behavioral correspondence between them.
\end{cvitems}

%\newpage

%\section{Past Research Activities}
%\cvheading{Privatizer in Parallelizing Compilers}{Jan. 2023 -- Aug. 2023}
%\cvdetail{Arcana Lab, Northwestern University \enspace | \enspace Advisor: Prof. Simone Campanoni}
%\begin{cvitems}
%  \item Helped design and implement an intraprocedural may-points-to analysis for \textbf{NOELLE}, and developed a privatizer that converts dynamic memory allocations and global variables into local variables while preserving program semantics.
%  \item Removed overly conservative loop-carried memory dependencies to enable further parallelization. Achieved roughly a \textbf{threefold speedup on eight cores} on a loop-intensive MiBench benchmark for which NOELLE previously achieved no speedup.
%\end{cvitems}
%\entrygap
%\cvheading{Sketch-based Image Synthesis}{Feb. 2020 -- Jun. 2020}
%\cvdetail{HCI Lab, Nanjing University \enspace | \enspace Advisor: Prof. Guihuan Feng}
%\begin{cvitems}
%  \item Compared semantic image synthesis algorithms and implemented spatially adaptive normalization. Built a pipeline that generates photorealistic landscapes from user-sketched segmentation masks, with control over semantics and style.
%\end{cvitems}
%
%\section{Selected Projects}
%\begin{cvitems}
%  \item \textbf{System F in Agda:} Collaboratively formalized System F and proved its type soundness.
%  \item \textbf{Lumi\`ere:} Began developing a functional language with two intermediate representations, tentatively named Fresnel and Poisson.
%  \item \textbf{FocusHour} (Feb. -- Apr. 2021): Built a productivity app with virtual trees, augmented reality, and visualizations of focus records and task completion. Received the highest score and top creativity award in a peer review involving more than 100 students.
%  \item \textbf{Orange'S} (Mar. -- Jun. 2018): Extended an experimental operating system on Bochs with Fibonacci computation, file system management, exception handling, and process scheduling.
%\end{cvitems}

\vsp

\section{Honors and Awards}
\textbf{National Scholarship for Postgraduates}, Nanjing University (top 1\%, Oct. 2021).\par
\textbf{First-class People's Scholarship} (top 5\%, Nov. 2019); \textbf{Third-class People's Scholarship} (top 20\%, Nov. 2017 and Nov. 2018), Nanjing University.\par
%\vspace{3pt}
%\textbf{Northwestern coursework (all A):} Programming Languages; Code Analysis and Transformation; Compiler Construction; Dynamics of Programming Languages; Proving Properties of Programs with Mechanized Logic; Advanced Topics in Compilers.

\newpage

\section{Work Experience}
\cvheading{Tencent Holdings Ltd.}{Jul. 2021 -- Jan. 2022}
\cvdetail{iOS Development Intern, Technology Engineering Group \enspace | \enspace Shenzhen, China}
\begin{cvitems}
  \item Developed an iOS multimedia messaging app using 5G connectivity; designed messaging card styles and interaction models, and refactored the user guide and network error detection logic.
\end{cvitems}

\vsp

\cvheading{Microsoft Software Technology Center Asia}{Jul. 2019 -- Oct. 2019}
\cvdetail{Software Engineering Intern, Speech Recognition Group \enspace | \enspace Suzhou, China}
\begin{cvitems}
  \item Automated speech recognition model training for conference scenarios using Microsoft AEther. Applied Gaussian blur and SG filtering to improve recognition accuracy in short conversations.
\end{cvitems}

\vsp



\section{Skills and Languages}
\textbf{Programming and formal methods:} Racket, Rocq/Coq, Redex, Agda, C++, Haskell, Python, Swift, LLVM IR.\par
\textbf{Tools:} Linux shell, Vim, \LaTeX, Git, HTCondor.\par
\textbf{Languages:} Chinese (native), English (proficient), Japanese (JLPT N1), French (beginner).

\end{document}
